% ===================================================================== Appendices A-H
% RE-ANALYSIS 2026
\titleformat{\chapter}[hang]{\sffamily\bfseries\LARGE}{Appendix \thechapter}{0.8em}{}[\vspace{2pt}\titlerule]
\chapter{Baseline Parameters}\label{app:params}
Table~\ref{tab:appA} lists the complete input set of the official baseline P00; the geometry is in Table~\ref{tab:geometry},
the operating point in Table~\ref{tab:operating} and the material data in Tables~\ref{tab:air} and~\ref{tab:inconel}. All inputs are class B
(selected with a stated rationale); no input is class A (taken from an original internship document).

%% generated by gen_tables.py - RE-ANALYSIS 2026
\begin{table}[htbp]
\centering
\footnotesize
\caption{Complete input set of the baseline P00}
\label{tab:appA}
\begin{tabular}{>{\raggedright\arraybackslash}p{0.34\linewidth}>{\raggedright\arraybackslash}p{0.56\linewidth}}
\toprule
Input & Value and basis \\
\midrule
Geometry & $D_i$ 20 / $D_o$ 40 / $t$ 10 / $L$ 600 mm (D-008) \\
Fluid & air, incompressible ideal gas; property tables 250--600 K (D-009, D-021) \\
Operating point & 300 K, 23.5 m/s, 0 Pa gauge at 101,325 Pa \\
Thermal boundary & 8000 W/m$^2$ uniform on the outer wall; adiabatic end faces (D-015) \\
Solid & Inconel 718; $\rho$ 8190 kg/m$^3$; $k(T)$, $c_p(T)$, $E(T)$, secant $\alpha(T)$, $S_y(T)$ (D-010, D-026, D-036) \\
Structural reference temperature & 300 K (D-041) \\
Supports & LC1 three-node determinate; LC2 = S1 (D-042, D-043) \\
CFD mesh & medium, 159,840 cells (D-034) \\
Structural mesh & B, 108,252 nodes / 23,400 SOLID186 (D-040, D-056) \\
Parametric ranges & $V$ $\pm$10\,\%, $q''$ $\pm$10\,\%, $t$ 8 / 12 mm with $Q$ held (D-058, D-060) \\
\bottomrule
\end{tabular}
\par\smallskip\parbox{\linewidth}{\footnotesize\raggedright Source: \texttt{MASTER\_ASSUMPTIONS.md} \S2; decision numbers refer to \texttt{PROJECT\_STATE.md}. All inputs are selected (class B) with a stated rationale.}
\end{table}


\chapter{Analytical Calculations}\label{app:analytical}
The Section 2 analytical baseline was computed by a self-testing Python script before any mesh existed (\texttt{baseline\_calculations.py},
18 sanity checks passed, energy balance closed to 0.0005\,\%). Its main results are listed in Table~\ref{tab:appB}; Figure~\ref{fig:anplots}
shows three of its plots. Worked examples of the key relations:

\begin{itemize}
\item Heat input: $Q = q''_o\,\pi D_o L = 8000 \times \pi \times 0.040 \times 0.600 = 603.19$ W.
\item Bulk temperature rise: $\Delta T_b = Q/(\dot m\,\bar c_p) = 603.19/(8.686\times10^{-3} \times 1008.5) = 68.85$ K, with $\bar c_p$ the mean
      specific heat over the bulk temperature range, giving $T_{out}$ = 368.85~K.
\item Inlet Reynolds number: $Re = G D_h/\mu = 27.650 \times 0.020/(1.846\times10^{-5}) = 29{,}957$.
\item Restrained stress: $\sigma_z = -E\alpha\Delta\bar T = -188.1\times10^{9} \times 13.72\times10^{-6} \times 254.7 = -657.2$ MPa, and free growth
      $\Delta L = \alpha\Delta\bar T L = 3.494\times10^{-3} \times 600 = 2.096$ mm.
\item Critical load from the FE eigenvalue: $P_{cr} = \lambda_1 N = 1.10805 \times 548{,}936.6 = 608.25$ kN; mean axial stress
      $N/A = 548{,}936.6/942.48 = 582.44$ MPa.
\end{itemize}

%% generated by gen_tables.py - RE-ANALYSIS 2026
\begin{table}[htbp]
\centering
\footnotesize
\caption{Section 2 analytical baseline (ANALYTICAL level, 1-D correlations; not a CFD or FE result)}
\label{tab:appB}
\begin{tabular}{lrl}
\toprule
Quantity & Value & Unit \\
\midrule
Mass flow rate & 8.686 & g/s \\
Reynolds number inlet / outlet & 29,957 / 25,548 & -- \\
Mean Darcy friction factor (Petukhov) & 0.02413 & -- \\
$\Delta p$ friction / acceleration / entry & 263.2 / 149.1 / 29.2 & Pa \\
$\Delta p$ CFD-comparable & 441.53 & Pa \\
Heat input $Q$ & 603.186 & W \\
Outlet bulk temperature & 368.85 & K \\
$Nu$ exit: Dittus--Boelter / Gnielinski / corrected & 66.79 / 61.87 / 49.58 & -- \\
$h$ exit, corrected & 77.83 & W/(m$^2$\,K) \\
Inner / outer wall temperature, exit & 574.43 / 581.71 & K \\
Through-wall $\Delta T$ (exit) & 7.285 & K \\
$R_{cond}/R_{conv}$ & 0.0342 & -- \\
Mean solid temperature (length-averaged mid-wall) & 554.74 & K \\
Free axial growth & 2.0964 & mm \\
LC1 peak von Mises (Timoshenko) & 16.313 & MPa \\
LC2 axial stress $-E\alpha\Delta T$ & $-$657.227 & MPa \\
$S_y$ at 281.6\,\textdegree C / LC2 utilisation & 1023.7 / 0.6420 & MPa / -- \\
LC2 / LC1 ratio & 40.3 & -- \\
Richardson / Brinkman number & $4.37\times10^{-5}$ / $1.82\times10^{-3}$ & -- \\
\bottomrule
\end{tabular}
\par\smallskip\parbox{\linewidth}{\footnotesize\raggedright Source: \texttt{02\_Engineering\_Calculations/baseline\_results.csv}; \texttt{MASTER\_PROJECT\_DATA.csv} rows M071--M080.}
\end{table}


\begin{figure}[htbp]
\centering
\begin{subfigure}[t]{0.34\linewidth}\centering\includegraphics[width=\linewidth]{01_axial_temperatures.png}\caption{Axial temperature development}\end{subfigure}\hfill
\begin{subfigure}[t]{0.31\linewidth}\centering\includegraphics[width=\linewidth]{03_radial_wall_temperature.png}\caption{Radial wall temperature}\end{subfigure}\hfill
\begin{subfigure}[t]{0.31\linewidth}\centering\includegraphics[width=\linewidth]{05_pressure_breakdown.png}\caption{Pressure-drop breakdown [Pa]}\end{subfigure}
\caption{Section 2 analytical baseline plots (ANALYTICAL level, produced before any CFD run).}
\label{fig:anplots}
\src{Source: \texttt{02\_Engineering\_Calculations/calculation\_plots/}; matplotlib plots of the analytical model.}
\end{figure}

\chapter{Mesh Tables}\label{app:mesh}
The CFD mesh family is defined in Tables~\ref{tab:meshfam} and~\ref{tab:meshq} and the structural mesh variants in Table~\ref{tab:smesh}.
On the coarse, medium and fine CFD meshes $y^+_{max}$ is 0.522, 0.585 and 0.653, always in the first slab, with 100\,\% of the wall at $y^+ \le 1$.

\chapter{CFD Monitor and Convergence Data}\label{app:conv}
The convergence criteria and final residuals are in Table~\ref{tab:convcrit} and the history in Figure~\ref{fig:conv}. The evaluation at
iteration 400 was not possible (too few second-order iterations), 500 was refused (scheme-switch transient), 600 met every criterion and
700 confirmed it. The fluid stayed within 299.998--553.42~K and the solid within 300.0--562.133~K throughout, inside both property tables,
with no limiter, divergence or reversed-flow message. All other CFD cases used the same staging and criteria.

\chapter{Structural Results}\label{app:struct}
The main static results are in Table~\ref{tab:lcres} and the first-yield assessment in Table~\ref{tab:yield}. Table~\ref{tab:appE} adds the
equilibrium and consistency checks. Stress singularities were checked on mesh B: averaged and unaveraged peaks differ by 0.30\,\% (LC1) and
0.004\,\% (LC2), and the energy-norm error is concentrated in the first 15~mm. An independent verification script re-derived the static results
from the raw MAPDL tables.

%% generated by gen_tables.py - RE-ANALYSIS 2026
\begin{table}[htbp]
\centering
\footnotesize
\caption{Additional structural results of the baseline}
\label{tab:appE}
\begin{tabular}{>{\raggedright\arraybackslash}p{0.55\linewidth}r}
\toprule
Quantity & Value \\
\midrule
LC1 reaction resultant [N] & $9.0\times10^{-7}$ \\
LC2 end reactions balanced to [N] & $1.8\times10^{-8}$ \\
LC2 reaction vs analytical restrained force of the same field & agreement to $10^{-5}$ \\
LC2 interior von Mises [MPa] & 587--591 \\
LC2 axial stress at mid-span, outer / bore [MPa] & $-$593.4 / $-$564.9 \\
LC2 maximum radial displacement [mm] & 0.0897 \\
LC1 free growth vs independent $\int\varepsilon_{th}\,dz$ & agreement to $1.3\times10^{-5}$ \\
LC2P peak von Mises [MPa] & 605.160918 \\
LC2 peak von Mises [MPa] (same precision) & 605.160873 \\
Energy-norm error LC1 / LC2 & 0.70\,\% / 0.012\,\% \\
\bottomrule
\end{tabular}
\par\smallskip\parbox{\linewidth}{\footnotesize\raggedright Source: \texttt{08\_Structural\_Analysis/STRUCTURAL\_RESULTS.md}, \texttt{PROJECT\_STATE.md} \S14.1, \texttt{FINAL\_ENGINEERING\_AUDIT.md} \S10.}
\end{table}


\chapter{Buckling Results}\label{app:buckling}
The hand estimates are in Table~\ref{tab:section}, the FE baseline results in Table~\ref{tab:buckling} and the support scenarios in
Table~\ref{tab:support}. Table~\ref{tab:appF} lists the first eigenvalue on every structural mesh. The workflow was benchmarked on a
constant-property toy tube: the Section 8A benchmark reproduced the closed-form Euler loads for the guided and no-sway end conditions, and
benchmark B03 gave a $\lambda_1$ between the plain and shear-corrected fixed--pinned values for the S3 remote-point formulation.

%% generated by gen_tables.py - RE-ANALYSIS 2026
\begin{table}[htbp]
\centering
\footnotesize
\caption{First eigenvalue on every structural mesh (S1 supports)}
\label{tab:appF}
\begin{tabular}{lrr}
\toprule
Mesh & Nodes & $\lambda_1$ (7 digits) \\
\midrule
XC & 24,171 & 1.1079876 \\
C & 50,652 & 1.1080342 \\
B (baseline) & 108,252 & 1.1080471 \\
FR & 127,080 & 1.1080404 \\
FA & 126,468 & 1.1080483 \\
FC & 126,294 & 1.1080410 \\
\bottomrule
\end{tabular}
\par\smallskip\parbox{\linewidth}{\footnotesize\raggedright Source: \texttt{PROJECT\_STATE.md} \S16.1 and \texttt{BUCKLING\_MESH\_COMPARISON.md}. Family monotonic (apparent order 4.9, GCI $5.6\times10^{-6}$); fine meshes within $6\times10^{-6}$ of B; $P_{cr}$ = 608.22--608.25 kN on all meshes.}
\end{table}


\chapter{Parametric Tables}\label{app:param}
The parametric results are in Tables~\ref{tab:pcfd}, \ref{tab:pstruct}, \ref{tab:pbk} and~\ref{tab:psens}; the changes from P00 are discussed in
Chapter~\ref{ch:param}, and the machine-readable tables are in the project folder \texttt{10\_Parametric\_Study}. The independent verification
\texttt{verify\_9B2.py} passed 45 of 46 checks; the single failure is the thickness check-mesh tolerance noted in Chapter~\ref{ch:unc},
reported rather than relaxed.

\chapter{Traceability}\label{app:trace}
\textbf{Status of the work.} The note on the analysis record at the front of the report applies to every value; only the engineer,
programme, organisation, period and tool set are documented internship facts.

\textbf{Numbers.} Every result value in the tables is read by \texttt{gen\_tables.py} from the Section 10A master dataset
\texttt{11\_Final\_Audit/MASTER\_PROJECT\_DATA.csv}: of its 182 rows, 124 are VERIFIED (recomputed from the raw solver or report file),
3 CROSS-CHECKED, 29 REPORTED (post-processing quantities from the section result files), 20 INPUT, 2 DECISION and 4 INTERPRETATION, and none
is a MISMATCH. Other values (solver settings, mesh controls, property tables) are quoted from the section documents named in each table note.
\texttt{qc\_10B.py} traces every number of the report to the master dataset or a project document (\texttt{Final\_Report/REPORT\_QC\_CHECKLIST.md}).

\textbf{Modelling levels and figures.} 581.7~K (analytical), 562.58~K (CFD medium) and 560.83~K (CFD fine), or $-$657.2~MPa (analytical
restrained stress) and 605.16~MPa (FE peak von Mises), belong to different modelling levels: they are different models or quantities, not
contradictions. All figures are hash-verified project files (\texttt{Figures/figure\_sources.json}); none is an ANSYS GUI screenshot.
